% Part VI -- Algorithms, Optimization, and Machine Learning
% Chapter 17: Numerical optimization with SciPy

\h{Numerical Optimization with SciPy}
\label{ch:17}

Optimization asks which feasible decision produces the best modeled outcome.
The answer is not created by a solver alone. A defensible workflow defines the
decision variables, objective function, constraints, units, and validation
tests before selecting an algorithm. This chapter begins with a transparent
grid search and then uses SciPy to refine a one-variable design recommendation.

\begin{NoteBox}
\textbf{Prerequisites.} You should understand functions, arrays, plotting, and
the algorithm contracts and efficiency ideas from \booklink{ch:16}{Chapter 16}.
The example is an educational model, not a building-performance standard.
\end{NoteBox}

\begin{tcolorbox}[title=Learning outcomes,colframe=orange,colback=white,arc=2mm]
After completing this chapter, you should be able to:
\begin{itemize}
    \item identify decision variables, an objective, bounds, and constraints;
    \item write an objective function that returns one scalar value;
    \item implement and interpret a finite grid search;
    \item use \c{scipy.optimize.minimize_scalar()} with bounded search;
    \item inspect the complete optimization result rather than only \c{result.x};
    \item verify feasibility and compare nearby candidate solutions;
    \item explain local minima, model sensitivity, and discretization error; and
    \item separate a computational optimum from a real-world recommendation.
\end{itemize}
\end{tcolorbox}

\hh{Formulate the model before calling a solver}

The running example chooses a window-to-wall ratio $r$. Larger windows may
improve daylight but can increase heating, cooling, or glare effects. We combine
those modeled consequences into one objective value, where lower is preferred,
and restrict $r$ to a feasible interval.

\bookfigure[0.98\textwidth]{Figures/Generated/ch17_optimization_model.pdf}
{An optimization workflow connects a decision variable, objective, feasible bounds, search algorithm, and independently verified result.}
{fig:ch17-model}

\begin{center}
\begin{tabular}{ll}
\textbf{Model component} & \textbf{Question} \\
Decision variable & What can the algorithm change? \\
Objective & What numerical quantity should be minimized or maximized? \\
Bounds and constraints & Which decisions are allowed? \\
Parameters and data & Which quantities are fixed for this run? \\
Validation & How will we check feasibility and nearby alternatives? \\
\end{tabular}
\end{center}

An objective is a model of priorities. Changing penalty weights can change the
optimum even when the solver is functioning perfectly.

\begin{ExplainBox}{A solver optimizes the stated model, not the real world}
The algorithm can search only the variables, objective, bounds, and constraints
it receives. Convergence is evidence about that mathematical search; it is not
evidence that the objective represents every design consequence, that the data
are accurate, or that the recommended decision is robust. Validate the model
with domain knowledge and compare nearby feasible alternatives before turning
a numerical optimum into a recommendation.
\end{ExplainBox}

\hh{Install SciPy once}

SciPy provides numerical optimization algorithms with a consistent result
interface.

\begin{verbatim}
python -m pip install scipy
\end{verbatim}

Import the specific function required by the problem.

\begin{lstlisting}[
    language=python,
    style=block,
    caption={Import a bounded scalar optimizer},
    label={c17_import},
    captionpos=b
]
from scipy.optimize import minimize_scalar
\end{lstlisting}

The official SciPy
\href{https://docs.scipy.org/doc/scipy/tutorial/optimize.html}{optimization guide}
describes scalar, multivariable, constrained, least-squares, root-finding, and
global methods. Choose an interface that matches the mathematical problem.

\hh{Write an objective that returns one scalar}

The objective below combines an energy component with penalties for insufficient
daylight and excessive glare. Each squared term changes smoothly except where a
penalty turns on or off.

\begin{lstlisting}[
    language=python,
    style=block,
    caption={Define a transparent educational objective},
    label={c17_objective},
    captionpos=b
]
def objective(window_ratio):
    energy = 80 + 160 * (window_ratio - 0.18) ** 2
    daylight_penalty = (
        450 * max(0.0, 0.32 - window_ratio) ** 2
    )
    glare_penalty = (
        600 * max(0.0, window_ratio - 0.55) ** 2
    )
    return energy + daylight_penalty + glare_penalty
\end{lstlisting}

\begin{SyntaxBox}{A scalar optimizer needs a scalar contract}
The function accepts one candidate ratio and returns one numeric score. The
\c{max(0.0, ...)} expressions make each penalty zero when its threshold is
satisfied. Squaring makes larger violations increasingly expensive. Returning
the sum gives the optimizer one quantity to compare.
\end{SyntaxBox}

\begin{TryBox}{Inspect the model before optimizing}
Evaluate the objective at ratios 0.10, 0.30, 0.50, and 0.70. Print the three
components separately. Explain which term dominates at each end of the range.
If the component values are surprising, revise the model before trusting a
solver.
\end{TryBox}

\hh{Begin with a finite grid search}

Grid search is easy to audit: generate feasible candidates, evaluate each one,
and retain the smallest value. Its limitation is resolution; the best point
between grid locations cannot be selected.

\begin{lstlisting}[
    language=python,
    style=block,
    caption={Search a finite set of feasible candidates},
    label={c17_grid_search},
    captionpos=b
]
import numpy as np

grid = np.linspace(0.10, 0.70, 13)
scores = np.array([objective(value) for value in grid])

best_index = scores.argmin()
best_ratio = grid[best_index]
best_score = scores[best_index]

print(best_ratio, best_score)
\end{lstlisting}

\begin{SyntaxBox}{Keep candidates and scores aligned}
\c{np.linspace()} includes both interval endpoints. The comprehension calls the
objective once per candidate. \c{argmin()} returns the \emph{index} of the
smallest score, not the score itself; that index retrieves the corresponding
ratio and objective value from aligned arrays.
\end{SyntaxBox}

The grid algorithm performs a linear scan of its candidate scores. More grid
points improve resolution but require more objective evaluations. In expensive
simulations, each evaluation may dominate the total runtime.

\hh{Refine the result with bounded scalar optimization}

SciPy's bounded method searches only inside the provided interval. It evaluates
the objective adaptively rather than at every equally spaced point.

\begin{lstlisting}[
    language=python,
    style=block,
    caption={Minimize one variable within explicit bounds},
    label={c17_minimize_scalar},
    captionpos=b
]
from scipy.optimize import minimize_scalar

result = minimize_scalar(
    objective,
    bounds=(0.10, 0.70),
    method="bounded",
)

print("Success:", result.success)
print("Recommended ratio:", result.x)
print("Objective value:", result.fun)
print("Function evaluations:", result.nfev)
print("Message:", result.message)
\end{lstlisting}

\begin{SyntaxBox}{Read the call and result object completely}
\begin{itemize}
    \item The first argument is the callable objective, not \c{objective()}.
          Passing the name lets the solver choose future arguments.
    \item \c{bounds=(low, high)} defines the feasible interval.
    \item \c{method="bounded"} selects an algorithm designed for a bounded
          one-variable problem.
    \item \c{result.x} is the reported decision and \c{result.fun} its score.
    \item \c{success}, \c{message}, and \c{nfev} describe the solver run.
\end{itemize}
Never treat \c{result.x} as valid when \c{success} is false or the point violates
the model's constraints.
\end{SyntaxBox}

\hh{Verify feasibility and local behavior}

Verification is a separate algorithm. Check the bounds, recompute the objective,
and compare nearby candidates.

\begin{lstlisting}[
    language=python,
    style=block,
    caption={Check the solver's recommendation independently},
    label={c17_verify},
    captionpos=b
]
if not result.success:
    raise RuntimeError(result.message)

if not 0.10 <= result.x <= 0.70:
    raise ValueError("Optimizer returned an infeasible ratio")

recomputed = objective(result.x)
if not abs(recomputed - result.fun) < 1e-8:
    raise ValueError("Reported and recomputed objectives differ")

epsilon = 0.01
left = objective(max(0.10, result.x - epsilon))
right = objective(min(0.70, result.x + epsilon))
print(left, result.fun, right)
\end{lstlisting}

The tolerance comparison accounts for floating-point representation. Nearby
checks do not prove a global minimum, but they can expose an incorrect result,
constraint mistake, or objective discontinuity.

\hh{Applied example: daylight and energy trade-off}
\label{sec:c17-application}

The complete runnable program is
\c{examples/c17_daylight_energy_optimization.py}. It calculates a grid-search
baseline, runs bounded optimization, validates the result, plots the objective
components, and exports the figure.

\begin{lstlisting}[
    language=python,
    style=block,
    caption={Compare grid and bounded optimization in one application},
    label={c17_application},
    captionpos=b
]
import numpy as np
from scipy.optimize import minimize_scalar


def objective(window_ratio):
    energy = 80 + 160 * (window_ratio - 0.18) ** 2
    daylight = 450 * max(0.0, 0.32 - window_ratio) ** 2
    glare = 600 * max(0.0, window_ratio - 0.55) ** 2
    return energy + daylight + glare


lower_bound = 0.10
upper_bound = 0.70

grid = np.linspace(lower_bound, upper_bound, 13)
grid_scores = np.array([objective(value) for value in grid])
grid_best = grid[grid_scores.argmin()]

result = minimize_scalar(
    objective,
    bounds=(lower_bound, upper_bound),
    method="bounded",
)

if not result.success:
    raise RuntimeError(result.message)
if not lower_bound <= result.x <= upper_bound:
    raise ValueError("Infeasible result")

print(f"Grid recommendation: {grid_best:.3f}")
print(f"Bounded recommendation: {result.x:.3f}")
print(f"Objective value: {result.fun:.2f}")
\end{lstlisting}

\bookfigure[0.88\textwidth]{Figures/Generated/ch17_window_optimization.pdf}
{The total objective combines energy and penalty components; grid candidates provide an auditable baseline while bounded optimization refines the reported window ratio.}
{fig:ch17-window-optimization}

\begin{ExplainBox}{Interpret the figure as model evidence}
The star identifies the lowest point reported by the bounded search, while open
circles show the coarser grid candidates. The optimum near 0.283 is not a
universal design rule. It is conditional on the equations, coefficients,
bounds, and omitted phenomena in this educational model.
\end{ExplainBox}

\hh{Extend to several decision variables carefully}

SciPy's \c{minimize()} accepts an initial vector and returns a vector solution.
The following syntax introduces two bounded variables without claiming that the
model is complete.

\begin{lstlisting}[
    language=python,
    style=block,
    caption={Minimize a two-variable objective with bounds},
    label={c17_multivariable},
    captionpos=b
]
from scipy.optimize import minimize

def two_variable_objective(decisions):
    window_ratio, shading_depth_m = decisions
    return (
        objective(window_ratio)
        + 25 * (shading_depth_m - 0.60) ** 2
        - 8 * window_ratio * shading_depth_m
    )

result_2d = minimize(
    two_variable_objective,
    x0=[0.30, 0.50],
    bounds=[(0.10, 0.70), (0.00, 1.20)],
    method="L-BFGS-B",
)
\end{lstlisting}

\c{x0} is the initial candidate vector. The objective unpacks that vector into
named decisions and still returns one scalar. Multivariable problems may have
several local minima, stronger interactions, and greater sensitivity to initial
conditions; visualize slices and try justified starting points.

\hh{Understand local minima and model sensitivity}

A local optimizer finds a point better than nearby candidates according to its
algorithm and tolerances. It may not find the best point across a complicated
domain. A strong workflow therefore:

\begin{enumerate}
    \item plots or samples the objective when dimensions permit;
    \item compares with a transparent baseline such as grid or random search;
    \item repeats from justified initial points for local multivariable methods;
    \item varies important coefficients and bounds; and
    \item reports how the recommendation changes.
\end{enumerate}

Sensitivity analysis asks whether the recommendation remains similar under
plausible model changes. A precise solution to an unstable model is not a robust
decision.

\hh{Common mistakes and useful diagnoses}

\begin{itemize}
    \item \textbf{The objective returns a list or array.} A scalar minimizer
          requires one numeric score per candidate.
    \item \textbf{Maximization is passed to a minimizer.} Minimize the negative
          utility or reformulate the objective and document the sign.
    \item \textbf{Bounds are omitted or use wrong units.} State feasibility and
          units before solving.
    \item \textbf{Only \c{result.x} is printed.} Inspect success, message,
          objective value, and evaluation count.
    \item \textbf{Solver output is treated as real-world truth.} Verify the
          model assumptions, data, and sensitivity.
\end{itemize}

\hh{Practice ladder}

\hhh{Level 0 -- Read}

Identify the decision variable, objective components, and feasible interval in
Listing~\ref{c17_objective} and Listing~\ref{c17_minimize_scalar}.

\hhh{Level 1 -- Modify}

Increase the daylight penalty coefficient from 450 to 700. Re-run the grid and
bounded searches and explain the direction of change.

\hhh{Level 2 -- Explain}

Explain why a 13-point grid can miss the bounded optimizer's recommendation and
why the bounded recommendation can still be misleading.

\hhh{Level 3 -- Build}

Write \c{grid_minimize(function, low, high, count)} that validates its arguments
and returns the best candidate and score.

\hhh{Level 4 -- Challenge}

Perform a sensitivity analysis over at least three daylight penalty weights.
Plot weight against recommended ratio and state whether the recommendation is
stable enough to communicate.

\hh{Practice solutions}

\textbf{Level 2:} The grid evaluates only fixed locations, so its resolution
limits the recommendation. The bounded search uses the written objective more
precisely, but an unrealistic equation or missing constraint makes its precise
answer irrelevant.

\textbf{Level 3}

\begin{lstlisting}[
    language=python,
    style=block,
    caption={One reusable grid minimizer},
    label={c17_solution_grid},
    captionpos=b
]
import numpy as np

def grid_minimize(function, low, high, count):
    if count < 2:
        raise ValueError("count must be at least 2")
    if low >= high:
        raise ValueError("low must be smaller than high")

    candidates = np.linspace(low, high, count)
    scores = np.array([function(value) for value in candidates])
    best_index = scores.argmin()
    return candidates[best_index], scores[best_index]
\end{lstlisting}

The function validates its search space, preserves alignment between candidates
and scores, and returns both the decision and evidence used to rank it.

\hh{Chapter checkpoint}

\begin{enumerate}
    \item What are the five main components of an optimization workflow?
    \item Why must the objective return one scalar?
    \item What trade-off controls grid-search resolution?
    \item What does \c{result.success} communicate?
    \item Why compare nearby candidates after solving?
    \item What is the difference between a local optimum and a global optimum?
    \item Why should penalty weights receive sensitivity analysis?
\end{enumerate}

\hh{Checkpoint answers}

\begin{enumerate}
    \item Decision variables, objective, constraints or bounds, search
          algorithm, and independent validation.
    \item The algorithm needs one comparable score for every candidate.
    \item More points improve resolution but require more evaluations.
    \item Whether the solver reports that its convergence conditions were met.
    \item To catch feasibility, objective, or convergence problems and inspect
          local behavior.
    \item A local optimum beats nearby candidates; a global optimum beats every
          feasible candidate.
    \item They encode priorities and can materially change the recommendation.
\end{enumerate}

\begin{tcolorbox}[title=Chapter summary,colframe=orange,colback=white,arc=2mm]
Optimization combines modeling and algorithms. Formulate a scalar objective and
feasible domain, compare a transparent baseline with an appropriate solver,
inspect the result object, and verify both the computation and model sensitivity.
\booklink{ch:18}{Chapter 18} applies the same disciplined separation of data,
algorithm, evaluation, and interpretation to machine learning.
\end{tcolorbox}
